% Part: lambda-calculus
% Chapter: lambda-definability
% Section: primitive-recursive-functions

\documentclass[../../../include/open-logic-section]{subfiles}

\begin{document}

\olfileid{lam}{ldf}{prf}
\olsection{بنسټيزې بازګشتي تابعې \usetoken{S}{lambda definable} دي} 

په ياد ولرئ چې بنسټيزې بازګشتي تابعې هغه دي چې له
بنسټيزو تابعو $\Zero$، $\Succ$ او $\Proj{n}{i}$ څخه د
ترکيب او بنسټيز بازګښت په وسيله تعريفېږي۔

\begin{lem}
  \ollabel{lem:basic}
  بنسټيزې بازګشتي تابعې $\Zero$، $\Succ$ او
  پروجکشن تابعې~$\Proj{n}{i}$ !!{lambda definable} دي۔
\end{lem}

\begin{proof}
دا تابعې د لاندې ترمونو په وسيله !!{lambda define} کېږي:
\begin{align*}
  \fn{Zero} & \ident \lambd[a][\lambd[fx][x]]\\
  \fn{Succ} & \ident \lambd[a][\lambd[fx][f (a f x)]] \\
  \fn{Proj}^n_i & \ident \lambd[x_0\dots x_{n-1}][x_i]
\end{align*}
\end{proof}

\begin{lem}
  \ollabel{lem:comp} فرض کړئ چې $k$-ځايه تابع $f$ او
  $n$-ځايه تابعې $g_0, \dots, g_{k-1}$ په ترتيب سره د
  $F$، $G_0$، \dots، او $G_{k-1}$ ترمونو په وسيله
  !!{lambda definable} دي، او $h$ له هغو څخه د ترکيب
  په وسيله تعريف شوې ده۔ بيا $h$ هم !!{lambda definable} ده۔
  \textbf{د سرچينې سمونونه (OLLAM-052، OLLAM-053):}
  په سرچينه کښې د وروستي استازي ترم کښته نښه او
  د پايلې تابع نوم د فرضونو له شمېر او ورپسې ثبوت
  سره نۀ سمون خوري؛ دلته دواړه له هماغو سره برابر دي۔
\end{lem}

\begin{proof}
  تابع $h$ د لاندې ترم په وسيله !!{lambda define} کېږي:
  \[
  H \ident \lambd[x_0 \dots x_{n-1}][F\, (G_0 x_0 \dots
    x_{n-1}) \dots (G_{k-1} x_0 \dots x_{n-1})]
  \]
  د دې خبرې کره کول د تمرين په توګه پرېږدو۔
\end{proof}

\begin{prob}
  د \olref[lam][ldf][prf]{lem:comp} ثبوت پوره کړئ؛
  وښيئ چې $H\num{n_0}\dots\num{n_{n-1}} \red \num{h(n_0, \dots,
    n_{n-1})}$۔
\end{prob}

پام وکړئ چې په \olref{lem:comp} کښې دا نۀ و شرط چې
$f$ او $g_0$، \dots،~$g_{k-1}$ دې بنسټيزې بازګشتي
وي؛ يوازې دا شرط و چې بشپړې او !!{lambda definable} وي۔

\begin{lem}\ollabel{lem:prim}
  فرض کړئ چې $f$ يوه $n$-ځايه تابع او~$g$ يوه
  $n+2$-ځايه تابع ده، او په ترتيب سره د $F$ او~$G$
  ترمونو په وسيله !!{lambda definable} دي۔ که تابع~$h$
  له $f$ او $g$ څخه د بنسټيز بازګښت په وسيله تعريف شوې
  وي، نو $h$ هم !!{lambda definable} ده۔
\end{lem}

\begin{proof}
  په ياد ولرئ چې $h$ داسې تعريفېږي:
  \begin{align*}
    h(x_1, \dots, x_n, 0) &= f(x_1, \dots, x_n)\\
    h(x_1, \dots, x_n, y+1) & = g(x_1, \dots, x_n, y, h(x_1, \dots, x_n, y)).
  \end{align*}
  \textbf{د سرچينې سمون (OLLAM-054):} د دويمې معادلې
  ښي لوري ته په سرچينه کښې د پايلې تابع بيا ليکل شوې،
  خو د فرض شوې ګامي تابع او لاندې ثبوت له مخې هلته
  هماغه ګامي تابع پکار ده۔ په نارسمي ډول، د $h$
  ارزښت د $y$ ګامونو په تېرېدو د ګامي تابع له مخکني
  ارزښت څخه ترلاسه کېږي، او پيلنی ارزښت
  $f(x_1, \dots, x_n)$ دے۔ د دې تکراري بڼې او د چرچ
  عددنښو د تکراروونکي تعريف تر منځ ورته‌والے شته۔

  د اسانۍ دپاره تعريف او ثبوت د يو اضافي آرګومېنټ~$x$
  په حالت کښې ورکوو۔ تابع $h$ د لاندې ترم په وسيله
  !!{lambda define} کېږي:
  \begin{align*}
    H \ident & \lambd[x][\lambd[y][\fn{Snd} (y
        D \tuple{\num{0}, F x})]]
    \intertext{چې پکې }
    D \ident & \lambd[p][\tuple{\fn{Succ} (\fn{Fst}\, p), (G x
          (\fn{Fst}\, p) (\fn{Snd}\, p))}]
  \end{align*}
  د تکرار ساتلے حالت يوه جوړه ده: لومړے غړے ئې
  د $y$-ګامي تکرار اوسنے شمېر او دويم ئې د~$h$
  اړوند ارزښت دے۔
  د تکرار په هر ګام کښې د $y$ او $h$ له نوو ارزښتونو
  يوه نوې جوړه جوړوو؛ له تکرار وروسته د جوړې دويم
  غړے اخلو، او هماغه د $h$ وروستنے ارزښت وي۔ اوس
  ثابتوو چې دا په رښتيا د بنسټيز بازګښت تمثيل دے۔

  غواړو ثابته کړو چې د هر $n$ او $m$ دپاره
  $H\,\num{n}\,\num{m} \red \num{h(n,m)}$۔ د دې دپاره
  لومړے ښيو چې که $D_n \ident \Subst{D}{\num{n}}{x}$
  وي، نو $D_n^m \tuple{\num{0}, F\, \num{n}} \red
  \tuple{\num{m}, \num{h(n, m)}}$۔ پر~$m$ استقرا کوو۔

  که $m=0$ وي، بايد $D_n^0 \tuple{\num{0}, F\, \num{n}} \red
  \tuple{\num{0}, \num{h(n, 0)}}$ وښيو۔ خو
  $D_n^0 \tuple{\num{0}, F\, \num{n}}$ هماغه
  $\tuple{\num{0}, F\, \num{n}}$ ده۔ څرنګه چې $F$
  تابع~$f$ !!{lambda define}s، دا جوړه
  $\tuple{\num{0}, \num{f(n)}}$ ته راکمېږي؛ او ځکه
  چې $f(n) = h(n, 0)$، دا هماغه
  $\tuple{\num{0}, \num{h(n,0)}}$ ده۔

  اوس فرض کړئ چې $D_n^m \tuple{\num{0}, F\, \num{n}} \red
  \tuple{\num{m}, \num{h(n, m)}}$۔ غواړو وښيو چې
  $D_n^{m+1} \tuple{\num{0}, F\, \num{n}} \red
  \tuple{\num{m+1}, \num{h(n, m+1)}}$۔
  \begin{align*}
    D_n^{m+1} \tuple{\num{0}, F\, \num{n}}
    & \ident D_n(D_n^{m} \tuple{\num{0}, F\, \num{n}})\\
    & \red D_n\,\tuple{\num{m}, \num{h(n, m)}} \text{\qquad (by IH)}\\
    & \ident (\lambd[p][
      \tuple{
        \fn{Succ} (\fn{Fst}\, p), (G \, \num{n} (\fn{Fst}\, p) (\fn{Snd}\, p))
    }]) \tuple{\num{m}, \num{h(n,m)}}\\
    & \redone
    \tuple{\fn{Succ} (\fn{Fst}\, \tuple{\num{m}, \num{h(n,m)}}),\\
      & \qquad (G \, \num{n} (\fn{Fst}\, \tuple{\num{m}, \num{h(n,m)}}) (\fn{Snd}\, \tuple{\num{m}, \num{h(n,m)}}))}\\
    & \red 
    \tuple{\fn{Succ}\, \num{m}, (G \, \num{n} \,\num{m} \, \num{h(n,m)})}\\
    & \red \tuple{\num{m+1}, \num{g(n, m, h(n, m))}}
  \end{align*}
  ځکه چې $g(n, m, h(n, m)) = h(n, m+1)$،
  غوښتل شوې پايله ترلاسه شوه۔

  په پاې کښې لاندې راکمېدنه وګورئ:
  \begin{align*}
    H\,\num n\,\num m &\ident \lambd[x][\lambd[y][\fn{Snd} (y
        (\lambd[p].\tuple{\fn{Succ} (\fn{Fst}\, p), (G\, x\,
          (\fn{Fst}\, p)\, (\fn{Snd}\, p))}) \tuple{\num{0}, F x})]]\\
    & \qquad\,\num n\, \num m\\
    & \red \fn{Snd} (\num m \,
    \underbrace{(\lambd[p].\tuple{\fn{Succ} (\fn{Fst}\, p), (G \,\num n\,
      (\fn{Fst}\, p) (\fn{Snd}\, p))})}_{D_n} \tuple{\num{0}, F \num n})\\
    & \ident \fn{Snd} (\num m \, D_n\, \tuple{\num{0}, F \num n})\\
    & \red \fn{Snd} \,(D_n^m  \tuple{\num{0}, F \num n})\\
    & \red \fn{Snd} \,\tuple{\num{m}, \num{h(n, m)}} \\
    & \red \num{h(n,m)}. 
  \end{align*}  
\end{proof}


\begin{prop}
  هره بنسټيزه بازګشتي تابع !!{lambda definable} ده۔
\end{prop}

\begin{proof}
  د \olref{lem:basic} له مخې ټولې بنسټيزې تابعې
  !!{lambda definable} دي، او د \olref{lem:comp}
  او \olref{lem:prim} له مخې !!{lambda definable}
  تابعې د ترکيب او بنسټيز بازګښت پر وړاندې تړلې دي۔
\end{proof}

\end{document}
